\section{Evaluation of Assisted Visual Analytics}
\label{sec:evaluation}
Evaluation tests whether a computational contribution helps scientists perform a task and what evidence supports that claim. Isenberg et al.'s review distinguishes goals including computational performance, user performance, and understanding analytical work.~\cite{EvalIsenberg2013} These outcomes are related but not interchangeable. A fast operation can produce an inappropriate comparison; a usable interface can expose a pattern that does not support a biological explanation. The task comparison in Section~\ref{sec:tasks} identifies the accomplishment to assess. The agency questions in Section~\ref{subsection:locusofagency} locate responsibility for proposal, scope, execution, and assessment, and the opportunities to inspect, change, or undo an operation.

The placements in Figure~\ref{fig:assistance-modality-map} establish mechanism--environment relationships, not comparative effectiveness. Each benefit therefore requires evaluation of the task and intervention actually claimed.

The system comparisons identify which contribution an evaluation needs to examine. For a configured pipeline, relevant questions include whether the chosen operations execute correctly and whether their outputs support the intended analysis. When a scientist teaches a classifier through labeled examples, evaluation must also examine how those examples affect subsequent assignments and whether erroneous propagation can be detected and corrected. A generated query requires checking the translation from the scientific request to data and operations, while comparison of supporting evidence requires inspection of provenance, coverage, and disagreement. These requirements follow from the analytical decisions supported or delegated, rather than from a shared interface or assistance label.

\subsection{Analytic Effectiveness and Scientific Utility}
\label{sec:evaluation-effectiveness}
Three checks are needed: whether an operation executes correctly, whether it implements the intended request, and whether its evidence warrants the scientific interpretation. Conversational generation makes the distinction visible: executable code can select the wrong cohort, and a correct cohort comparison may still leave confounding unresolved. For adaptive assistance, prediction accuracy and the benefit of recommending the predicted action require separate evaluation.

The baseline determines which contribution can be isolated. A matched comparison with and without guidance can assess the intervention; changing the display, representation, and interaction simultaneously assesses a system-level combination. Poco et al.'s projection study motivates capable desktop alternatives when evaluating spatial approaches.~\cite{EvalPoco2011} Participant expertise, data complexity, duration, and task definitions must accompany performance results. A prescribed cluster-identification task assesses that bounded accomplishment; the cluster's biological interpretation requires domain evidence.

The RetainVis--MediSyn comparison in Section~\ref{sec:tasks-comparison} illustrates the evidential boundary. Predictive-model measurements and performance on simplified evidence-selection tasks answer different questions. Computational benchmarks establish output properties; controlled studies test bounded human-performance claims; expert cases identify domain relevance; sustained use can reveal effects on an evolving workflow. None should be substituted for another merely because all evaluate the same software.

Operational uptake provides a further kind of evidence. Shyr et al. report REDCap generative-AI use in 958 Vanderbilt projects, including 754 using qualitative summarization and 215 using writing assistance; feature counts overlap.~\cite{Shyr2026REDCapAI} These totals include exploratory and operational use. They establish deployment and use, while summary fidelity, verification effort, and effects on research conclusions require separate assessment. For assisted analysis, a useful next test is whether researchers can identify omissions or unsupported interpretations in generated summaries of their data.

The expanded studies reinforce these distinctions. Interactive AI annotation reports 192 tasks performed by three radiologists, with improved annotation quality for two participants but not the third and concerns about unlocalized model changes.~\cite{ExpandedE24} Five biologists assessing the Space-Time Hypercube found utility for wave detection but mixed results for division and asynchrony tasks.~\cite{ExpandedE22} BioNet-XR's 30-person usability study tested desktop use; immersive versions were presented through videos.~\cite{ExpandedE26} Liver navigation reports phantom registration performance rather than clinical outcomes, and ASCRIBE-XR demonstrates implemented workflows without a user-benefit study.~\cite{ExpandedE15,ExpandedE28} These results cannot be combined into a single claim of analytical effectiveness.

Assessment of metastasis detection for the Augmented Reality Microscope exposes a further complication: agreement with a recorded label need not imply a correct result. Reinspection against immunostains identified discrepancies in the ground truth itself.~\cite{ExpansionX014} Model accuracy, spatial presentation, and assistance to pathological decisions therefore require distinct checks; a detector assessment is not sufficient evidence of improved decision-making through the microscope.

\subsection{Representational Fidelity and Uncertainty}
\label{sec:evaluation-fidelity}
Filtering, aggregation, layout, and embedding change which relationships are available for inspection. Fidelity concerns preservation of the properties needed for the task, rather than a universal preference for maximum detail. Holten's edge bundling reduces clutter while creating ambiguity in following individual edges.~\cite{EvalHolten2006} CMGV's distinction between underlying and visible graphs supports complexity-management operations; integrity and performance tests assess those operations rather than scientists' understanding.~\cite{EvalZafar2026}

Computational and interpretive fidelity require separate evidence. An aggregate may correctly represent a population mean while concealing a small subgroup important to the question. Evaluation can test preservation of relevant structure and whether analysts recognize when omitted detail matters. Biological uncertainty, representational distortion, and uncertainty about intent also need separation. A well-registered spatial highlight may still identify the wrong object; a correct selection may still feed an inappropriate transformation. Combining all three errors into one accuracy measure obscures the correction required.

Useful assessments therefore connect the displayed result to its source observations and active transformations, testing whether scientists can detect a consequential omission, inspect alternatives, and revise an interpretation. The relevant properties depend on the task: orientation requires stable correspondence across scales, while comparison requires a defensible basis for the contrast.

\subsection{Provenance and Recovery}
\label{sec:evaluation-provenance}
Provenance is valuable when scientists can use it to reconstruct and question a result. VisTrails records evolving scientific workflows and supports returning to earlier versions and comparing alternatives.~\cite{EvalFreire2006} This provides a basis for recovery, not proof that a user will notice an error. Conversational explanations add another distinction: an account of what was intended is not evidence of what actually executed. DeepVIS exposes intermediate generation steps and permits targeted refinement, but favorable assessments of inspectability do not establish reliable recovery in biological analysis.~\cite{EvalShuai2025}

Adaptive traces should identify the context and model state used for an intervention, the alternatives shown, and acceptance, rejection, or modification. Subsequent behavior may have been induced by the intervention rather than independently expressing a preference. Reconstruction can require software and model versions, stochastic seeds, parameters, and learned-state updates as well as prompts and outputs.

Recovery should be tested beyond undoing the last action. If an incorrect filter feeds a comparison and generated summary, correcting the filter must also expose stale dependent results. Relevant measures include error detection, localization of the responsible operation, successful restoration or revision, correction time, and persistence of downstream errors. In shared work, the trace also needs to distinguish who proposed, executed, and accepted a change. These checks turn recorded history into usable accountability.

\subsection{Analyst Control, Calibration, and Exploration Bias}
\label{sec:evaluation-control}
Control includes setting assistance scope before an action, accepting or modifying it during execution, and recovering afterward. ProactiveVA users requested finer intervention during agent activity; rapidly changing tasks and noisy traces also complicated intent inference.~\cite{EvalZhao2025} A useful initial suggestion therefore does not establish that later operations remain appropriate.

Calibrated reliance means accepting sound assistance while questioning incorrect or unsupported output. In PhenoFlow, lengthy explanations could conceal errors, while shortening them could remove necessary information; code and data-field inspection provided complementary checks.~\cite{EvalKim2025} Evaluation should measure acceptance of correct suggestions, detection and correction of incorrect ones, and recognition that further evidence is needed. Trust ratings or acceptance frequency alone cannot establish calibration.

Ha et al. compare user models for interaction prediction and exploration-bias detection, separating these capabilities from the benefit of subsequent guidance.~\cite{EvalHa2023} The feedback loop in Section~\ref{sec:assistance-adaptive} motivates following multiple interventions: a recommendation alters the next selection and hence the evidence used to recommend again. Studies can record exposure to alternatives, introduce known intent mismatches, and examine behavior after goals change. Exploration breadth should be interpreted with recognition of contradictory evidence and quality of justification, not rewarded as novelty for its own sake. Expertise in biology, experience with the interface, and familiarity with the model should be reported separately.

\subsection{Workload and Embodied Cost}
\label{sec:evaluation-workload}
Assistance can transfer work from execution to specification, inspection, and repair. Timing only query generation or counting clicks misses the effort required to reach a trustworthy result. Adaptive timing also matters: a correct suggestion may interrupt a demanding comparison. Reorientation, rejected prompts, and explaining changed goals can offset the saved interaction.

Computational delays should be separated by operation. Graphia compares loading, layout, and rendering but notes different layout stopping conditions, limiting direct runtime comparison.~\cite{EvalFreeman2022} Likewise, Yang et al.'s room-sized and zoomable scatterplot conditions reveal task-dependent navigation costs, without isolating an AI contribution.~\cite{EvalYang2021} Assisted viewpoint changes might reduce travel yet increase the effort to recover context.

Selection accuracy, correction actions, task time, and reported effort should therefore be assessed together. Spatial workflows additionally require reporting input devices, seated or standing use, session duration, discomfort, breaks, and withdrawals. Fewer movements may reflect useful computation or inadequate inspection; shorter completion time may coexist with greater mental effort. A practical evaluation follows the entire analytical episode, including waiting, supervision, and recovery.

\subsection{Collaboration and Common Ground}
\label{sec:evaluation-collaboration}
Collaborators need to understand the purpose, state, and evidential basis of shared work. Visibility is only one component. FIESTA's exploratory study of ten groups of three identified movement between personal and joint work and perceived awareness benefits in shared VR.~\cite{EvalLee2021} It did not test autonomous group oversight. In \emph{Talk to the Wall}, ten pairs used speech and touch in document analysis; fourteen participants reported awareness benefits from hearing a partner's commands, while distraction was also a concern.~\cite{EvalMolinaLeon2025} A causal improvement in task speed or accuracy was not demonstrated in that study.

The consequential question is whose intention governs a shared intervention. Following the shared-selection example in Section~\ref{sec:modality-walls}, evaluation should distinguish a private preview from an authorized group change, test whether affected participants notice its scope, and assess whether they can challenge or reverse it. A request from the most active speaker is not evidence of collective agreement.

Relevant outcomes include clarification effort, awareness of consequential changes, resolution of conflicting operations, restoration of prior state, and participants' ability to explain a joint conclusion. Equal command counts are not inherently desirable when roles differ; access to decisions and the ability to contribute or contest evidence matter more. Agreement aids coordination but does not validate a biological claim.

\subsection{Cross-Study Evidence and Limits of Comparability}
\label{sec:evaluation-comparison}
Table~\ref{tab:evaluation-evidence} compares selected evaluations by their intervention, study context, evidence, and limits. It preserves distinctions among computational measurement, controlled task performance, participant feedback, and observations of use. Spatial-language requirements elicited through Wizard-of-Oz methods and implemented scene-operation architectures answer different questions, as discussed in Section~\ref{sec:modality-vr}; neither automatically establishes scientific benefit.

Quantitative comparison is justified when the task, outcome definition, and baseline align. Response time requires equivalent timed work; accuracy requires a comparable target. Where studies are heterogeneous, qualitative comparison can still establish which mechanism was tested and where its evidence stops. Heterogeneity does not justify discarding compatible measurements, and an unevaluated capability is not evidence of failure.

General visualization studies can inform interaction choices without proving biological utility. Domain cases can establish relevance without isolating causal benefit. Across these forms of evidence, the recurring need is to connect the computational contribution to analyst use, error correction, and the scientific judgment that follows. Section~\ref{sec:challenges} develops research directions from those gaps rather than treating greater autonomy or immersion as a goal in itself.
\begin{table*}[t]

\centering

\caption{Selected evaluation examples and the limits of their supported claims. Participant feedback, computational measurements, and observed behavior are identified separately. The comparison concerns the evidence used in this section, rather than every evaluation reported in each paper.}

\label{tab:evaluation-evidence}

\footnotesize

\setlength{\tabcolsep}{4pt}

\renewcommand{\arraystretch}{1.12}

\begin{tabular}{@{}>{\raggedright\arraybackslash}p{.17\textwidth}>{\raggedright\arraybackslash}p{.21\textwidth}>{\raggedright\arraybackslash}p{.24\textwidth}>{\raggedright\arraybackslash}p{.32\textwidth}@{}}

\toprule

\textbf{Study and analytic context} & \textbf{Intervention and comparison} & \textbf{Evidence and study context} & \textbf{Supported claim and limit} \\

\midrule

Graphia~\cite{EvalFreeman2022}: desktop graph exploration & Graph loading, layout, and rendering compared with Gephi and Cytoscape & Computational comparison using three graphs, including a gene-expression graph; hardware and software versions reported & Demonstrates performance characteristics under the tested conditions. Different layout stopping conditions limit runtime comparison; rendering speed does not establish reduced cognitive workload. \\

\addlinespace[3pt]

Yang et al.~\cite{EvalYang2021}: navigation and comparison in VR scatterplots & Room-sized versus zoomable views, each with and without an overview & Controlled comparison of four navigation conditions; task time, accuracy, and navigation analysis & Supports task-dependent navigation choices. It does not isolate an AI intervention or establish biological utility. \\

\addlinespace[3pt]

Ha et al.~\cite{EvalHa2023}: modeling data exploration & Eight user-modeling algorithms compared across four interaction datasets & Computational evaluation of interaction prediction, bias detection, and complexity & Compares user-model performance. Successful prediction or detection does not establish that a subsequent intervention improves exploration. \\

\addlinespace[3pt]

DeepVIS~\cite{EvalShuai2025}: conversational visualization generation & Step-wise generation and editable intermediate decisions; computational baselines and interface evaluation & Generation benchmarks and a 32-participant study; participant assessments of understanding and refinement & Provides distinct evidence about output generation and perceived inspectability. These results do not establish reliable error recovery in biological workflows. \\

\addlinespace[3pt]

PhenoFlow~\cite{EvalKim2025}: clinical cohort exploration & LLM-assisted data wrangling with code, explanation, and data inspection & Design collaboration with two neurologists and clinical case studies; observations during iterative refinement & Grounds inspection requirements in domain practice. Explanation-related difficulties are design evidence, not a controlled estimate of calibration improvement. \\

\addlinespace[3pt]

ProactiveVA~\cite{EvalZhao2025}: proactive assistance during event analysis & Assistance informed by interaction history; comparison with a formative study & Ten-user evaluation with 15-minute exploration, interaction records, ratings, and interviews & Reports use of proactive assistance and identifies steering limitations. A formative-study comparison and short sessions do not establish sustained scientific benefit. \\

\addlinespace[3pt]

FIESTA~\cite{EvalLee2021}: collaborative analysis in a shared VR workspace & Configurations supporting individual and joint work & Exploratory study with ten groups of three using housing data; behavior observations and participant feedback & Identifies coordination practices and perceived awareness benefits. Timing differences were not strictly controlled; autonomous group oversight was not tested. \\

\addlinespace[3pt]

Talk to the Wall~\cite{EvalMolinaLeon2025}: collaborative analysis on a large display & Speech and touch used within a shared workspace & Ten-pair study; interaction records and participant reports & Shows how speech can make actions available to collaborators while raising distraction concerns. Positive reports do not establish a causal gain in team accuracy or autonomous oversight. \\

\bottomrule

\end{tabular}

\end{table*}
